% Einheit 3 – Exponentialfunktion aufstellen und auswerten
\setcounter{aufgabe}{27}
\hypertarget{e3}{}\einheitenkopf*{Einheit 3 von 5 · Exponentialfunktion aufstellen und auswerten}
\zweigzeile{Hier lernst du, eine Gleichung für ein Wachstum aufzustellen und damit Werte und Zeitpunkte zu berechnen \,·\, neu in diesem Jahr \,·\, P10 \,·\, baut auf: Einheit 2, Potenz mit dem Taschenrechner, Zinseszins}

\begin{aufgabe}{Ich erkenne, welches Jahr der Schritt null ist. Kreise im Text das Startjahr ein und unterstreiche das gesuchte Jahr. Du musst nichts ausrechnen.}
\begin{teile}
\teil 2022 hatte ein Verein 180 Mitglieder. Jedes Jahr kommen $7\,\%$ dazu. Wie viele Mitglieder hat er 2029?
\teil Im Jahr 2021 wurden 90 Störche gezählt; ihre Zahl wächst jährlich um $12\,\%$. Wie viele sind es 2026?
\teil Eine Wohnung kostete 2019 genau $850\,€$ Miete, danach jedes Jahr $2\,\%$ mehr. Wie hoch ist die Miete 2027?
\teil Wie viel ist ein Auto 2030 wert, wenn es 2023 neu $28\,000\,€$ kostete und jährlich $12\,\%$ an Wert verliert?
\teil Schreibe in der Zeitleiste für c) unter das Startjahr $t = 0$ und kreuze das gesuchte Jahr an.
\end{teile}

\sachtabelle{lccccccccc}{Jahr & 2019 & 2020 & 2021 & 2022 & 2023 & 2024 & 2025 & 2026 & 2027}{Schritt & & & & & & & & & }

\end{aufgabe}

\begin{aufgabe}{Ich kann eine Exponentialfunktion aufstellen. Unterstreiche im Text den Anfangswert und den Prozentsatz. Gib den Anfangswert $N_0$ und den Wachstumsfaktor $q$ an.}
\beispiel{\text{Ein Teich hat 360 Fische, ihre Zahl wächst jedes Jahr um } 12\,\% &\;\rightarrow\; N_0 = 360, \quad q = 1{,}12}
\begin{geruest}
\gz{a}{Ein Sparbuch mit $700\,€$ wird jährlich mit $2\,\%$ verzinst.}{\feld{N_0}\feld{q}}
\gz{b}{Eine Algenfläche von $12\,\text{m}^2$ wächst täglich um $15\,\%$.}{\feld{N_0}\feld{q}}
\gz{c}{Eine Stadt mit $48\,000$ Einwohnern wächst jährlich um $2{,}5\,\%$.}{\feld{N_0}\feld{q}}
\gz{d}{Eine Kultur mit $900$ Bakterien wächst stündlich um $35\,\%$.}{\feld{N_0}\feld{q}}
\end{geruest}
Stelle die Gleichung $N(t) = N_0 \cdot q^t$ auf. Dabei ist $N(t)$ der Bestand nach $t$ Zeitschritten (Jahren, Tagen, Stunden).
\begin{teile}
\teil für b) \quad \feldl{N(t)}
\teil Ein Gletscher bedeckt $40\,\text{km}^2$ und verliert jedes Jahr $2\,\%$ seiner Fläche. \quad \feldl{N(t)}
\end{teile}
\end{aufgabe}

\begin{aufgabe}{Ich kann eine Exponentialfunktion aufstellen – weiter.}
\begin{teile}
\teil Der Wert einer Maschine von $50\,000\,€$ sinkt jedes Jahr um $12\,\%$. In den Gleichungen ist $y$ der Wert in € und $x$ die Zeit in Jahren. Kreuze die passende Gleichung an. \\
\kreuz{$y = 50\,000 \cdot 0{,}12^x$} \\
\kreuz{$y = 50\,000 \cdot 1{,}12^x$} \\
\kreuz{$y = 50\,000 \cdot 0{,}88^x$} \\
\kreuz{$y = 50\,000 - 0{,}12\,x$}
\teil Die Gleichung $N(t) = 2\,500 \cdot 1{,}07^t$ beschreibt die Zahl der Besucher eines Museums $t$ Jahre nach der Eröffnung. Trage ein, was die Bestandteile bedeuten.
\end{teile}
\sachtabelle{ll}{Bestandteil & Bedeutung}{$2\,500$ & \langzelle \\ $1{,}07$ & \langzelle \\ $t$ & \langzelle}
\begin{teile}
\teil Erkläre, warum im Wachstumsfaktor $1{,}07$ die $1$ für $100\,\%$ steht.
\teil Die Stromkosten eines Haushalts betragen im Jahr 2023 genau $1\,150\,€$. Sie steigen jedes Jahr um $2{,}8\,\%$. Stelle eine Gleichung auf, mit der man die Stromkosten $t$ Jahre nach 2023 berechnen kann. (P10 2026)
\end{teile}
\end{aufgabe}

\begin{aufgabe}{Ich kann mit einer Exponentialfunktion Werte berechnen. Berechne den Wert Schritt für Schritt, indem du mehrmals mit $q$ malnimmst.}
\beispiel{N(t) = 250 \cdot 1{,}1^t: \quad N(2) &= 250 \cdot 1{,}1 \cdot 1{,}1 = 302{,}5}
\begin{teilezwei}
\tz $N(t) = 500 \cdot 2^t$: \ \feld{N(3)} & \tz $N(t) = 90 \cdot 1{,}1^t$: \ \feld{N(2)} \\
\tz $N(t) = 2\,000 \cdot 1{,}05^t$: \feld{N(2)} & \tz $N(t) = 50 \cdot 3^t$: \ \feld{N(2)} \\
\end{teilezwei}
Berechne mit der Potenztaste in einem Schritt.
\begin{teile}
\teil $N(t) = 2\,000 \cdot 1{,}05^t$: \quad \feld{N(3)}
\teil $N(t) = 3\,400 \cdot 1{,}02^t$ beschreibt die Einwohnerzahl eines Dorfes, $t$ in Jahren ab 2022. Berechne die Einwohnerzahl im Jahr 2030. \quad \leerfeld
\teil $N(t) = 1\,250 \cdot 1{,}035^t$ beschreibt ein Guthaben in €, $t$ in Jahren. Berechne das Guthaben nach $7$ Jahren und runde auf Cent. \quad \leerfeld[€]
\teil Ein Händler behauptet: „Eine Uhr für $2\,400\,€$ verliert jedes Jahr $7\,\%$ an Wert. Nach $5$ Jahren ist sie weniger als $1\,700\,€$ wert.“ Prüfe die Behauptung.
\teil Eine Stadt hat $18\,000$ Einwohner und wächst jedes Jahr um $2\,\%$. Studie A rechnet mit der Exponentialfunktion. Studie B rechnet: „12 Jahre mal $2\,\%$ sind $24\,\%$, also $22\,320$ Einwohner.“ Berechne die Vorhersage von Studie A für 12 Jahre und vergleiche.
\end{teile}
\end{aufgabe}

\begin{aufgabe}{Ich kann mit einer Exponentialfunktion Werte berechnen – weiter. Rechne Schritt für Schritt, bis die Schwelle überschritten ist. Gib an, nach wie vielen Schritten das passiert und welcher Wert dann erreicht ist.}
\begin{teile}
\teil $N(t) = 750 \cdot 1{,}35^t$ beschreibt die Zahl der Pflanzen in einem Beet, $t$ in Wochen. Wann sind es erstmals mehr als $2\,000$?
\teil Auf einem See schwimmen $900\,\text{kg}$ Algen. Durch Pflege werden es jeden Tag $15\,\%$ weniger. Nach wie vielen Tagen sind es erstmals weniger als $300\,\text{kg}$?
\teil Die Zahl der Nutzer einer App wird beschrieben durch $N(t) = 3\,500 \cdot 1{,}25^t$, $t$ in Monaten. Die Entwickler behaupten, nach einem Jahr hätten sie mehr als $50\,000$ Nutzer. Prüfe die Behauptung. (P10 2025)
\end{teile}
\end{aufgabe}

\begin{aufgabe}{Ich finde den Fehler beim Berechnen mit einer Exponentialfunktion. Ole soll für $N(t) = 600 \cdot 1{,}05^t$ den Wert nach 8 Jahren berechnen. Er rechnet so:}
\rechnung{N(8) &= 600 \cdot 1{,}05 \cdot 8 = 5\,040}
\begin{teile}
\teil Finde den Fehler, benenne ihn und korrigiere die Rechnung.
\teil Berechne selbst für $N(t) = 750 \cdot 1{,}1^t$ den Wert nach $6$ Jahren. \quad \leerfeld
\end{teile}
Ella soll für $N(t) = 1\,800 \cdot 1{,}02^t$, $t$ in Jahren ab 2025, den Wert im Jahr 2031 berechnen. Sie rechnet so:
\rechnung{N &= 1\,800 \cdot 1{,}02^{2031}}
\begin{teile}
\teil Finde den Fehler, benenne ihn und korrigiere die Rechnung.
\teil Berechne selbst für $N(t) = 2\,600 \cdot 1{,}015^t$, $t$ in Jahren ab 2023, den Wert im Jahr 2030. \quad \leerfeld
\end{teile}
\end{aufgabe}

\begin{aufgabe}{Ich kann begründen, warum bei exponentiellem Wachstum eine Potenz entsteht. Entscheide ohne Rechnung und begründe.}
\begin{teile}
\teil Warum steht in $N(t) = 500 \cdot 1{,}1^t$ das $t$ im Exponenten und nicht als Faktor vor der $1{,}1$?
\teil Julia sagt: „Wächst etwas 12 Jahre lang jedes Jahr um $2\,\%$, dann ist es danach um mehr als $24\,\%$ gewachsen.“ Hat sie recht?
\end{teile}
\end{aufgabe}

\begin{aufgabe}{Ich kann mit einer Exponentialfunktion eine Entscheidung treffen. Ein Patient nimmt eine Tablette mit $250\,\text{mg}$ Wirkstoff. Sein Körper baut jede Stunde $18\,\%$ des Wirkstoffs ab. Eine zweite Tablette darf er erst nehmen, wenn weniger als $90\,\text{mg}$ im Körper sind.}
\begin{teile}
\teil Stelle eine Gleichung für die Wirkstoffmenge nach $t$ Stunden auf.
\teil Wie viel Wirkstoff ist nach $3$ Stunden noch im Körper? Runde auf eine Nachkommastelle.
\teil Nach wie vielen vollen Stunden darf er die zweite Tablette nehmen?
\end{teile}
\end{aufgabe}
